\documentclass[UTF8]{ctexart}
\usepackage{xeCJK}
\usepackage{geometry}
\geometry{left=2cm,right=2cm,top=2.5cm,bottom=2.5cm}
\setlength{\headheight}{12.7pt}

\usepackage{titlesec}
\titleformat{\section}
  {\normalfont\large\bfseries}
  {\thesection}
  {1em}
  {}
\titlespacing*{\section}
  {0pt}
  {3.5ex plus 1ex minus .2ex}
  {2.3ex plus .2ex}

\usepackage{amsmath}
\usepackage{amssymb}
\usepackage{amsthm}
\usepackage{enumitem}
\setlist[enumerate]{itemsep=0pt,topsep=0pt,parsep=0pt,leftmargin=0pt,itemindent=2.5em}
\setlist[itemize]{itemsep=0pt,topsep=0pt,parsep=0pt,leftmargin=0pt,itemindent=2.5em}
\usepackage{fancyhdr}
\pagestyle{fancy}
\fancyhf{}
\usepackage{graphicx}
\usepackage{hyperref}
\usepackage{indentfirst}
\usepackage{lmodern}


\begin{document}
\setlength{\abovedisplayskip}{1pt}
\setlength{\belowdisplayskip}{1pt}

\section{基数}

\hypertarget{sec:定义1.1}{}
\noindent\textbf{定义1.1 (基数)}

给定\href{../../04 序理论/02 序数/01 序数#nameddest=sec:定义1.1}{序数}
$\kappa$, 如果
\[
    \forall\alpha<\kappa:\quad\alpha\prec\kappa,
\]
就称$\kappa$是一个\textbf{基数}.

\vspace{10pt}

\hypertarget{sec:定理1.1}{}
\noindent\textbf{定理1.1}

自然数和自然数集$\omega$都是基数.

\noindent{\kaishu 证明.}

\begin{enumerate}
    \item 自然数$n$都是有限集, 即$\forall n<\omega:n\prec\omega$,
    \item 任取自然数$n$,
    由\href{../../03 自然数/03 有限集/01 有限集#nameddest=sec:定理1.2}{鸽笼原理}得
    $\forall m<n:m\prec n$. \hfill $\qedsymbol$
\end{enumerate}

\vspace{10pt}

\hypertarget{sec:定义1.2}{}
\noindent\textbf{定义1.2 (集合的势)}

给定集合$x$, 存在唯一的基数$\kappa$使得$x\approx\kappa$,
称之为$x$的\textbf{势}, 记作$|x|$或$\#x$.

\noindent{\kaishu 证明.}

\noindent{\kaishu 存在性.}

考虑$x$上所有良序关系的集合
$L=\{r\subset x\times x\mid\,r\text{\,是\,}x\text{\,上的良序\,}\}$,
然后收集它们的序型
\[
    x^{\mathrm{E}}=\{\,\mathrm{Type}(x,r)\,\mid r\in L\}.
\]

根据\href{../../04 序理论/02 序数/05 良序原理#nameddest=sec:定理1.1}{良序原理},
$L$非空, 所以$x^{\mathrm{E}}$非空.
令$\kappa=\min x^{\mathrm{E}}$, 显然$\kappa\approx x$.

对任意的$\alpha<\kappa$, 假设$\alpha\approx\kappa\approx x$,
存在双射$f:\alpha\to x$, 把$\alpha$上的良序复制到$x$上, 记作$\leqslant$, 即有
\[
    f:(\alpha,\in)\cong(x,<)\quad\implies\quad
    \alpha=\mathrm{Type}(x,\leqslant)\quad\implies\quad
    \alpha\in x^{\mathrm{E}},
\]
矛盾. 所以$\alpha\prec\kappa$, 即$\kappa$是基数.

\noindent{\kaishu 唯一性.}

如果基数$\kappa,\lambda$都与$x$等势, 假设$\kappa<\lambda$则$\kappa\prec\lambda$,
矛盾. 同理$\lambda<\kappa$也矛盾, 所以$\kappa=\lambda$. \hfill $\qedsymbol$

\vspace{10pt}

\hypertarget{sec:定理1.2}{}
\noindent\textbf{定理1.2}

任给集合$a,b$ :
\begin{enumerate}
    \item $a\approx b$当且仅当$|a|=|b|$,
    \item $a\prec b$当且仅当$|a|<|b|$,
    \item $a\preccurlyeq b$当且仅当$|a|\leqslant|b|$.
\end{enumerate}

\vspace{10pt}

\hypertarget{sec:定义1.3}{}
\noindent\textbf{定义1.3 (基数类)}

记全体基数组成的类为\textbf{基数类}, 记作$\mathbf{Cn}$.
基数类是\href{../../01 ZFC 公理/02 类#nameddest=sec:真类}{真类}.

\noindent{\kaishu 验证.}

假设存在集合$\mathrm{Cn}$使得
$\forall\kappa:\kappa\in\mathrm{Cn}\leftrightarrow\kappa\text{\,是基数\,}$. 那么:
\[
    \left|\,\mathcal{P}(\sup\mathrm{Cn})\,\right|\in\mathrm{Cn}
    \quad\implies\quad
    \left|\,\mathcal{P}(\sup\mathrm{Cn})\,\right|\leqslant\sup\mathrm{Cn}
    \quad\implies\quad
    \mathcal{P}(\sup\mathrm{Cn})\preccurlyeq\sup\mathrm{Cn},
\]
矛盾. \hfill $\qedsymbol$





\newpage

\section{基数的构造}

Hartogs数和取上确界两种操作可以得到更大的基数.

\vspace{10pt}

\hypertarget{sec:定义2.1}{}
\noindent\textbf{定义2.1 (Hartogs数)}

\vspace{3pt}
任给集合$x$, 定义$\displaystyle
x^{\mathrm{h}}\overset{\mathrm{def}}{=}\bigcup_{y\subset x}y^{\mathrm{E}}$,
称之为$x$的\textbf{Hartogs数}.

\vspace{10pt}

\hypertarget{sec:定理2.1}{}
\noindent\textbf{定理2.1}

任给集合$x$, 有
\[
    \forall\alpha:\quad\alpha\in x^{\mathrm{h}}\leftrightarrow
    \alpha\in\mathbf{On}\land\alpha\preccurlyeq x.
\]

\noindent{\kaishu 证明.}

\begin{enumerate}
    \item 如果$\alpha\in x^{\mathrm{h}}$,
    则存在$y\subset x$使得$\alpha\in y^{\mathrm{E}}$, 从而
    $\alpha\approx y\preccurlyeq x$.
    \item 如果序数$\alpha\preccurlyeq x$, 那么有单射$f:\alpha\to x$,
    即有双射$f:\alpha\to\mathrm{Ran}\,f$, 所以
    \begin{equation*}
        \alpha\approx\mathrm{Ran}\,f\quad\implies\quad
        \alpha\in(\mathrm{Ran}\,f)\,^{\mathrm{E}}\quad\implies\quad
        \alpha\in x^{\mathrm{h}}.
    \tag*{\qedsymbol}\end{equation*}
\end{enumerate}

\vspace{10pt}

\hypertarget{sec:定理2.2}{}
\noindent\textbf{定理2.2}

任给集合$x$, $x^{\mathrm{h}}$是大于$|x|$的最小基数.

\noindent{\kaishu 证明.}

第一, 证明$x^{\mathrm{h}}$是序数. 如果$\beta\in\alpha\in x^{\mathrm{h}}$, 则
\[
    \beta\preccurlyeq\alpha\land\alpha\preccurlyeq x
    \quad\implies\quad\beta\preccurlyeq x
    \quad\implies\quad\beta\in x^{\mathrm{h}},
\]
即$x^{\mathrm{h}}$是传递集, 所以是序数.

第二, 证明$x^{\mathrm{h}}$是基数. 首先假设$x^{\mathrm{h}}\preccurlyeq x$
则$x^{\mathrm{h}}\in x^{\mathrm{h}}$矛盾, 所以$x^{\mathrm{h}}\succ x$.
然后对任意的$\alpha<x^{\mathrm{h}}$有
\[
    \alpha\preccurlyeq x\quad\implies\quad\alpha\prec x^{\mathrm{h}},
\]
所以$x^{\mathrm{h}}$是基数.

第三, 由$x\prec x^{\mathrm{h}}$得$|x|<x^{\mathrm{h}}$.
另外, 如果基数$\kappa>|x|$, 此时假设$\kappa<x^{\mathrm{h}}$则
\[
    \kappa\preccurlyeq x\quad\implies\quad\kappa\leqslant|x|,
\]
矛盾. 所以$\kappa\geqslant x^{\mathrm{h}}$,
即$x^{\mathrm{h}}$是大于$|x|$的最小基数. \hfill $\qedsymbol$

\vspace{15pt}

\hypertarget{sec:定理2.3}{}
\noindent\textbf{定理2.3}

任给非空集合$K$, 如果$K$的元素都是基数, 那么$\sup K$也是基数.

\noindent{\kaishu 证明.}

对任意的$\displaystyle\alpha<\sup K=\bigcup K$, 存在$\kappa\in K$使得
\[
    \alpha<\kappa\quad\implies\quad\alpha\prec\kappa
    \quad\implies\quad\alpha\prec\sup K,
\]
所以$\sup K$是基数. \hfill $\qedsymbol$

\end{document}